\section{讨论}

本节回到引言中提出的三个研究问题，并结合第~\ref{sec:results} 节证据作答。

\subsection{RQ1：为什么表示连续性最弱的架构反而最常获胜？}

CKA--性能悖论，即 GraphResGNN 在深层拥有最低 CKA（0.942）却赢得最多组合（12/24），而 NodeResGNN 拥有最高 CKA（0.968）却只赢得 3 次，具有明确的结构性解释。GraphResGNN 串联三个图条件块，每个块都接收前一块的池化图嵌入作为条件输入。该设计并不追求保留表示，而是利用逐步累积的全局上下文，在每个阶段刻意重塑表示。在当前 benchmark 中，这一策略在 DD、ENZYMES、AIDS 和 Mutagenicity 上更强，但并非在所有地方都最佳。

NodeResGNN 更忠实地保留表示，但也存在代价：如果表示在层间变化过小，模型可能无法充分适应任务特定结构。交叉残差架构在一定程度上缓解了这种张力，因为它为单分支可能压制掉的信息提供了另一条路径，却不会像 GraphResGNN 那样强制进行激进重塑。这也是 GraphCrossGNN 虽然很少拿到第一，但经常保持在前列的原因。

\textbf{回答：} 在本文设定中，积极的表示重塑可以是有益的，但其收益依赖于数据集。更高的连续性不足以保证更好的准确率，而交叉残差交换更应被理解为一种稳定化替代方案，而不是证明“连续性应当总是降低”。

\subsection{RQ1 \& RQ2：不同复用拓扑何时占优？}

该 benchmark 支持的是分裂型优势区间，而不是单一赢家。在面向论文主线的 biological comparison 中，NodeResGNN 在 PROTEINS 上最强，而 GraphResGNN 在 DD 上最强。交叉残差变体在 DD 上最有说服力：在算子级 winner count 中，NodeCrossGNN 赢下四个 DD 算子设定中的三个，尽管在单一 GCNConv 主表中，GraphResGNN 仍是最优 DD 模型。这个区别很重要，因为它表明 DD 上的交叉残差优势在跨算子层面是稳的，但又不能被简化为某一次单模型 sweep。ENZYMES 也指向 GraphResGNN 更强，但由于所有模型在该数据集上都表现较弱，它更适合被解释为压力测试，而不是与主证据并列。

在图级摘要更有用的任务上，GraphResGNN 是更强的默认选择；在最优残差基线已能匹配或超过最优交叉残差变体的数据集上，交叉交换应被视为有竞争力的备选，而不是默认首选。图~\ref{fig:cross_advantage_heatmap} 对这一边界做了明确说明：最优交叉变体通常优于 PlainGNN，但很少优于最优残差模型。补充数据集有助于检验这一边界并非蛋白质核心独有，但它们不应反过来重写本文的主 biological argument。

\textbf{回答：} 在当前 benchmark 中，图级残差传播可作为最强的默认家族；当任务更接近 PROTEINS 时，节点级残差复用仍是更强选择；而交叉残差交换是一种选择性方案，尤其值得在类似 DD 的蛋白质图和跨算子比较中测试。

\subsection{RQ3：路由策略是否在架构选择之外仍然重要？}

第~\ref{sec:results} 节中的消融表明答案是肯定的，而且这种作用依赖目标。第一，gate ablation 显示交叉残差架构对 gate 配置较为稳健，可学习与固定 gate 在偏向交叉的目标上各赢三次，这与其稳定的排序画像一致。第二，路由模式确实重要：稀疏路由（阈值 0.05）在四个代表性目标中赢下三个，其中两个涉及交叉残差架构。其机制层面的协同在于，稀疏性能够在残差信号跨深度累积噪声前先抑制弱通道，而跨分支路径则保留真正互补的信息。过强稀疏（0.10）会损害图级残差目标，这再次说明其作用依赖阈值，而不是普适定律。

\textbf{回答：} 路由策略是一个有意义的次级设计轴。轻度稀疏路由通常是交叉残差交换的更佳搭档，而可学习稠密路由则仍在最强的图级残差目标上保持优势。

\subsection{局限性}

这些结论的适用范围受到若干限制。首先，统计效能有限：在仅有五个折的前提下，只有较大效应才容易达到 $p < 0.05$。其次，CKA--性能悖论目前仅在单一数据集--算子组合（PROTEINS, GCNConv）上被系统记录，其能否推广仍待验证。第三，这些 benchmark 数据集属于中等规模生物分子图，在更大规模或不同类型图上，权衡关系可能改变。第四，路由分析只覆盖四个目标，足以揭示模式，但还不足以支撑普适性结论。

\subsection{未来方向}

最重要的扩展方向，是在更大 benchmark 和更丰富输入上继续检验 CKA--性能悖论。如果该现象能够泛化，它将改变我们评价深层 GNN 的方式，把关注点从层间稳定性指标转向表示变化是否具有结构性。稀疏路由与交叉残差的协同关系，也应在超过当前四个目标的更广泛设定中继续测试。
